\documentclass[11pt]{article}
\usepackage[a4paper,margin=18mm]{geometry}
\usepackage[no-math]{fontspec}
\setmainfont{Latin Modern Roman}
\usepackage{amsmath,amssymb,amsthm,xfrac,xspace,hyperref,graphicx,comment}
\excludecomment{DefTralics}

\providecommand{\bibliofont}{\small}
\newcommand{\ie}{i.e.,\xspace}
\newcommand{\eg}{e.g.,\xspace}
\newcommand{\etc}{etc.\xspace}
\newcommand{\cf}{cf.\xspace}
\newcommand{\resp}{resp.\xspace}
\newcommand{\smaller}{\small}
\newcommand{\Subsection}{\subsection}
\newcommand{\Subsubsection}{\subsubsection}
\newcommand{\skpt}{\smallskip}
\emergencystretch=3em
\numberwithin{equation}{section}\providecommand{\loccit}{\textit{loc.\ cit.}}
\input{author-preamble.tex}
\begin{document}
\begin{center}{\Large Stable item 2040}\end{center}
\paragraph{Source and attribution.}
Carlos E. Arreche, Thomas Dreyfus, Julien Roques, \emph{Differential transcendence criteria for second-order linear difference equations and elliptic hypergeometric functions}, Journal de l’École polytechnique — Mathématiques 8 (2021), publisher version, DOI 10.5802/jep.143. \href{https://jep.centre-mersenne.org/articles/10.5802/jep.143/}{Publisher article}. Authors' text: \href{https://creativecommons.org/licenses/by/4.0/}{CC BY 4.0}.
\paragraph{Editorial problem boundary.}
Prove only the field-extension instance of Theorem 3.5 below, under the original characteristic-zero, commuting automorphism/derivation, algebraically closed constants, property (P), absence of a base-field Riccati solution and absence of a nonzero constant-coefficient telescoper. Every nonzero solution of equation (3.1) in a difference-differential field extension is differentially transcendental. The exact Lemma 3.4 reduction and full imprimitive/scalar-group contradiction are retained. The printed theorem display \(a(y)\) is unchanged; equation (3.1), the companion system and proof give the author\'s \(a\sigma(y)\) notation. General nonreduced algebras and later elliptic applications are outside this selected claim.
\paragraph{Difficulty (editorial): H2.}
Standard Picard--Vessiot theory and the cited Galois classification only reduce the equation to primitive and imprimitive alternatives. The author must transfer a nonzero differentially algebraic solution to the parametrised Picard--Vessiot setting and exclude the imprimitive branch through a change-of-basis determinant and the Riccati/telescoper obstructions. In the primitive branch, the scalar Galois action and differential-polynomial evaluation yield a separate contradiction, with every nonzero-solution step retained in the selected field-extension scope.
\paragraph{Editorial transcription note.}
Original author language, mathematics and ordering are preserved. Publisher front matter is replaced by this independent article wrapper. Bibliography is recovered from matching cached publisher HTML, including its TeX formulas. No new proof or mathematical repair is supplied. Surrounding original results are retained for conservative context and are not extra selected corpus claims.
\clearpage
\input{author-body.tex}
\begin{thebibliography}{99}
\bibitem{arreche2017computation} C. E. Arreche - “Computation of the difference-differential Galois group and differential relations among solutions for a second-order linear difference equation” , Commun. Contemp. Math. 19 (2017) no. 6, article ID 1650056, 42 pages
\bibitem{arreche2017galois} C. E. Arreche \& M. F. Singer - “Galois groups for integrable and projectively integrable linear difference equations” , J. Algebra 480 (2017), p. 423-449
\bibitem{DHR} T. Dreyfus, C. Hardouin \& J. Roques - “Hypertranscendence of solutions of Mahler equations” , J. Eur. Math. Soc. (JEMS) 20 (2018) no. 9, p. 2209-2238
\bibitem{dreyfus2016functional} T. Dreyfus, C. Hardouin \& J. Roques - “Functional relations of solutions of $ q $-difference equations” , Math. Z. (2021), doi:10.1007/s00209-020-02669-4
\bibitem{dreyfus2018nature} T. Dreyfus, C. Hardouin, J. Roques \& M. F. Singer - “On the nature of the generating series of walks in the quarter plane” , Invent. Math. 213 (2018) no. 1, p. 139-203
\bibitem{dreyfus2017walks} T. Dreyfus, C. Hardouin, J. Roques \& M. F. Singer - “Walks in the quarter plane: genus zero case” , J. Combin. Theory Ser. A 174 (2020), p. 105251, 25
\bibitem{dreyfus2015galois} T. Dreyfus \& J. Roques - “Galois groups of difference equations of order two on elliptic curves” , SIGMA Symmetry Integrability Geom. Methods Appl. 11 (2015), article ID 003, 23 pages
\bibitem{dreyfus2017differential} T. Dreyfus \& K. Raschel - “Differential transcendence \& algebraicity criteria for the series counting weighted quadrant walks” , Publ. Math. Besançon (2019) no. 1, p. 41-80
\bibitem{Hen97} P. A. Hendriks - “An algorithm for computing a standard form for second-order linear $q$-difference equations” , J. Pure Appl. Algebra 117/118 (1997), p. 331-352, Algorithms for algebra (Eindhoven, 1996)
\bibitem{Hen98} P. A. Hendriks - “An algorithm determining the difference Galois group of second order linear difference equations” , J. Symbolic Comput. 26 (1998) no. 4, p. 445-461
\bibitem{HS} C. Hardouin \& M. F. Singer - “Differential Galois theory of linear difference equations” , Math. Ann. 342 (2008) no. 2, p. 333-377, Erratum: Ibid. 350 (2011), no. 1, p. 243–244
\bibitem{holder1886ueber} O. Hölder - “Ueber die Eigenschaft der Gammafunction keiner algebraischen Differentialgleichung zu genügen” , Math. Ann. 28 (1886), p. 1-13
\bibitem{Kol73} E. R. Kolchin - Differential algebra and algebraic groups , Pure and Applied Math. , vol. 54 , Academic Press, New York-London, 1973
\bibitem{Kol74} E. R. Kolchin - “Constrained extensions of differential fields” , Adv. Math. 12 (1974), p. 141-170
\bibitem{magnus2009elliptic} A. P. Magnus - “Elliptic hypergeometric solutions to elliptic difference equations” , SIGMA Symmetry Integrability Geom. Methods Appl. 5 (2009), article ID 038, 12 pages
\bibitem{Mumford} D. Mumford - Tata lectures on theta. I , Modern Birkhäuser Classics , Birkhäuser Boston, Inc., Boston, MA, 2007, Reprint of the 1983 edition
\bibitem{rains2010transformations} E. M. Rains - “Transformations of elliptic hypergeometric integrals” , Ann. of Math. (2) 171 (2010) no. 1, p. 169-243
\bibitem{Ro11} J. Roques - “Generalized basic hypergeometric equations” , Invent. Math. 184 (2011) no. 3, p. 499-528
\bibitem{roques2018algebraic} J. Roques - “On the algebraic relations between Mahler functions” , Trans. Amer. Math. Soc. 370 (2018) no. 1, p. 321-355
\bibitem{rosengren2002elliptic} H. Rosengren \& M. J. Schlosser - “Multidimensional matrix inversions and elliptic hypergeometric series on root systems” , SIGMA Symmetry Integrability Geom. Methods Appl. 16 (2020), article ID 088, 21 pages
\bibitem{spiridonov2016elliptic} V. P. Spiridonov - “Elliptic hypergeometric functions” , 2016
\bibitem{van2007hyperbolic} F. J. van de Bult et al. - Hyperbolic hypergeometric functions , University of Amsterdam, Amsterdam Netherlands, 2007
\bibitem{fokko2009basic} F. J. van de Bult \& E. M. Rains - “Basic hypergeometric functions as limits of elliptic hypergeometric functions” , SIGMA Symmetry Integrability Geom. Methods Appl. 5 (2009), article ID 059, 31 pages
\bibitem{VdPS97} M. van der Put \& M. F. Singer - Galois theory of difference equations , Lect. Notes in Math. , vol. 1666 , Springer-Verlag, Berlin, 1997
\end{thebibliography}
\end{document}
